\section{The complete fourth-order coefficient}
\label{sec:fourth-order}

We now fix the Wick convention $\epsilon U=1$. The factor
$(\epsilon\pi\nu)^2=1/(4096\pi^8)$ will be left explicit.
At order $\hbar^4$ there are tree graphs, marked two-loop
theta graphs, and bare four-loop graphs. Odd geometric loop
numbers vanish, irrespective of the counterterm orders at
the vertices.

\subsection{Trees and marked theta graphs}

Use the four matrix units and momenta of
\eqref{eq:all-order-extraction-denominator}. The differential
of the unscaled fourth-order cubic evaluates to $16$.
On these inputs the two tree contractions are
\[
 \{D,[p_0p_1]^3\tr(X_0\{X_1,X_2\})\}=10,\qquad
 \frac12\{[p_0p_1]^2f,[p_0p_1]^2f\}=-4.
\]
For the marked theta graphs the following entries include
the sum over allowed marked vertices and the graph
automorphism factors, with the displayed vertex coupling
and the common factor $\pi\nu$ removed:
\begin{center}
\renewcommand{\arraystretch}{1.7}
\begin{tabular}{@{}lcc@{}}
\toprule
Theta graph & One $[p_0p_1]^2f$ vertex & Two $D$ vertices\\
\midrule
$(1,3,3)$ & $-\dfrac{499N^2}{2835}-\dfrac{256}{105}$
  & $\dfrac{431N^2}{11340}+\dfrac{718}{2835}$\\
$(2,2,3)$ & $-\dfrac{34}{27}$ & $\dfrac5{63}$\\
\bottomrule
\end{tabular}
\end{center}
Every contraction uses the complete traceless inverse pairing,
including its $-1/N$ term. Substituting
\eqref{eq:first-two-counterterms} and \eqref{eq:i3}, adding the
trees, and dividing by $16$ gives the following contribution
to the primitive coefficient:
\begin{equation}
 (\epsilon\pi\nu)^2
       \frac{21221N^4+85304N^2+192480}{3628800}.
 \label{eq:four-tree-marked}
\end{equation}

\subsection{The fifty bare graphs}

A tight cubic four-loop graph has ten vertices, thirteen
internal edges, and four external fields. The internal degree
of every external vertex is two: an internal degree-one vertex
would violate the tight bound on the remaining induced
subgraph. A cut vertex is also excluded by applying the tight
bound to the two subgraphs that meet there. Suppressing the
four bivalent vertices therefore leaves a biconnected cubic
multigraph on six vertices.

A loop at a core vertex is incompatible with this
biconnectedness. Three parallel edges would isolate two
vertices. The remaining cores are enumerated by their simple
graphs and the matchings that specify doubled edges.
Distributing the four subdivisions, dividing by the core
automorphisms, and checking the tight inequality on every
induced subset leaves exactly fifty isomorphism classes.
The four surviving core types contain $3,8,28,11$ classes.
The oriented edges, subdivision paths, and full automorphism
orders are recorded individually in the accompanying
\texttt{four-loop-graphs.json} data of \cite{RepoHigher}.

The integrated kernels are homogeneous quartics in the three
momentum determinants. On the fixed cyclic matrix-unit input,
the four core sums, divided by $\pi^2$ before the common heat
normalization, have the following coefficients:
\begin{center}
\small
\renewcommand{\arraystretch}{2.1}
\setlength{\tabcolsep}{4pt}
\begin{tabular}{@{}lccc@{}}
\toprule
Graph indices & Coefficient of $1$ & Coefficient of $N^2$
  & Coefficient of $N^4$\\
\midrule
$0$--$2$ & $0$ & $0$ & $\frac1{10500}$\\
$3$--$10$
 & $\frac{101}{945}-\frac{1241\pi^2}{45360}$
 & $\frac{73903}{1020600}-\frac{37\pi^2}{3402}$
 & $\frac{34259}{20412000}-\frac{11\pi^2}{829440}$\\
$11$--$38$
 & $\frac{9796575\pi^2-96791482}{12247200}$
 & $\frac{88921725\pi^2-884922356}{734832000}$
 & $\frac{144949856-14698125\pi^2}{653184000}$\\
$39$--$49$
 & $\frac{93253402-9461505\pi^2}{12247200}$
 & $\frac{734213669-74453325\pi^2}{183708000}$
 & $0$\\
\bottomrule
\end{tabular}
\end{center}
Their exact sum is
\begin{equation}
 -\frac{172}{945}
 +N^2\left(\frac{139229}{48600}-\frac{35701\pi^2}{120960}\right)
 +N^4\left(\frac{760981}{3402000}
                       -\frac{130727\pi^2}{5806080}\right).
 \label{eq:four-bare-sum}
\end{equation}
In this convention the common four-loop heat factor in the
cochain equation is
\[
 \frac1{(8\pi^{5/2})^4}
       =\frac{(\epsilon\pi\nu)^2}{\pi^2}.
\]
The contribution of \eqref{eq:four-bare-sum} to the primitive
is therefore $(\epsilon\pi\nu)^2/16$ times that expression.

\begin{theorem}[Fourth-order counterterm]
\label{thm:i4}
In the stated heat-kernel scheme and Wick convention, the
fourth-order Taylor cochain is
\begin{equation}
\boxed{\begin{aligned}
 i_4={}&\frac{1}{4096\pi^8}
 \left[
   \frac1{24}
   +N^2\left(\frac{1102559}{5443200}
                       -\frac{35701\pi^2}{1935360}\right)\right.\\[-1mm]
 &\left.\hspace{23mm}
   +N^4\left(\frac{4216}{212625}
                       -\frac{130727\pi^2}{92897280}\right)
 \right][p_0p_1]^4\tr(X_0[X_1,X_2]).
\end{aligned}}
\label{eq:i4}
\end{equation}
The entire local order-four equation holds for arbitrary
traceless colors and conserved momenta, and specializes to
every $N\ge2$. The same $i_4$ applies to both choices of $a$.
\end{theorem}

\begin{proof}
Add \eqref{eq:four-tree-marked} to one sixteenth of
\eqref{eq:four-bare-sum}. The constant part is
\[
 \frac{192480}{3628800}-\frac{172}{16\cdot945}
       =\frac1{24};
\]
the $N^2$ and $N^4$ parts give exactly the two displayed
parentheses. Theorem~\ref{thm:all-orders-traceless} proves
closedness and uniqueness before this evaluation is used.
Thus the single evaluation determines the coefficient.
There is also a separate full-cochain check: after reconstructing
the color contractions and all differentiated tree and theta
terms, the rational part and the $\pi^2$ part each agree with
$di_4$ in all $300$ nonzero trace and momentum coordinates,
including their entire Laurent polynomials in $N$.
\end{proof}

\begin{corollary}
The coefficient of
$(\epsilon\pi\nu)^2[p_0p_1]^4f$ in $i_4$ is positive and
transcendental for every integer $N\ge2$.
\end{corollary}
\begin{proof}
The strict inequality $\pi^2<10$ gives
\begin{align*}
 \frac{1102559}{5443200}
          -\frac{35701\pi^2}{1935360}
    &>\frac{787747}{43545600}>0,\\
 \frac{4216}{212625}
          -\frac{130727\pi^2}{92897280}
    &>\frac{2864617}{497664000}>0.
\end{align*}
At each positive integer $N$, the coefficient is a rational
number minus
\[
 \left(\frac{35701N^2}{1935360}
              +\frac{130727N^4}{92897280}\right)\pi^2 ,
\]
whose rational multiplier of $\pi^2$ is nonzero.
Transcendence of $\pi$ then proves the assertion.
\end{proof}

The coefficient is attached to this renormalization scheme
and this choice of local representative. Its transcendence
does not identify a scheme-independent observable.
The full integrated kernels are inherited exact input;
the present verification independently reconstructs their
complete cochain assembly, rather than repeating every
four-loop master integration.
